\minitopic{完整示范：从“知之为知之”到认识谦逊}以《论语》“知之为知之，不知为不知”与当代认识谦逊作比较，第一步不是宣布二者相同，而是分别重建问题。原句处在问答、学习和言说规范中，至少涉及主体不把未知冒充已知；当代认识谦逊可能讨论证据限度、能力边界、对反证开放或避免过度自信。二者共享“正确标记自己的认识位置”这一桥梁问题，但不必共享知识定义、德性结构或社会角色\parencite[2.17]{analects2003}。

第二步寻找反向材料。若只收集赞成自知之明的句子，对应几乎必然成功。研究者还要问：文本是否把承认无知与求教、身份秩序或道德修养连接？当代理论是否允许一个谦逊主体仍然断言并行动？一方若强调学习关系，另一方若主要评价个体置信校准，差异就会改变理论。比较结论应写成：二者在“避免伪装知道”这一功能上相近，但对无知如何被克服、谁有资格纠正以及认识目标是什么，给出不同结构。

第三步才进入规范建构。可以提出一种关系性的认识谦逊：主体不仅准确报告自己的把握程度，还应知道何时寻找更好来源、怎样让对方能够纠正自己，以及权威关系何时妨碍承认无知。这个方案受到古典材料和当代认识论共同启发，却是研究者的新主张。把三步分开，比较既不沦为格言装饰，也不要求历史文本预先拥有现代术语\parencite{norden2017}。

\minitopic{竞争解释不是附录，而是内部解释的一部分}古典文本通常没有唯一无争议的现代释义。同一句话可以被理解为语义判断、修养劝告、政治言说规范或特定角色的回应。比较者应至少列出一种能改变现代对应的竞争解释，并说明自己依据哪些篇章、注疏、用词或整体论证作选择。若两种解释都与证据相容，比较结论就应使用条件句，而不是把争议隐藏在译文中。

现代对话理论也必须接受同样要求。“语境主义”“德性认识论”或“知识规范”各自包含内部版本。选择其中最方便匹配的弱版本，会制造虚假接近。合格比较要让双方都有具体作者、问题和论证，不能一边是精确文本，另一边只是教科书标签。证据责任的对称性，是双向批评能够真正发生的前提。

\minitopic{跨文化案例实验需要四道控制}跨文化Gettier实验已经在多种语言和群体中检验“有根据的真信念却不算知识”的判断\parencite{machery2017}，但直译和统计票数远远不够。第一是语言控制：目标语言中的“知道”是否覆盖相同对象和语用强度。第二是情境控制：支票、谷仓或司法程序等背景是否为各组熟悉。第三是理解控制：受试者是否把握信念为真的偶然路径，而不是被故事长度或人物道德吸引。第四是测量控制：二元选择、程度评分和解释性回答是否指向同一判断。

高质量研究可以先访谈少量参与者，寻找他们实际依据的特征；再共同设计在各语境中功能相当的案例；通过回译检查字面差异，同时保留无法对齐之处；最后报告组内分布和不确定区间，而不是把每个群体写成单一“文化直觉”。即使差异稳定，它首先支持的是关于该任务中判断分布的主张，不能直接证明古典传统造成差异，更不能单独决定知识的规范本质。

\begin{center}
\begin{tabularx}{\textwidth}{p{0.18\textwidth}XX}
\toprule
研究层次 & 可以支持的结论 & 不能直接推出的结论\\
\midrule
文本解释 & 某段材料在特定版本中的问题与论证 & 整个传统或当代群体都持同一立场\\
结构比较 & 两套理论在指定问题轴上的相似与差异 & 双方概念完全同一或存在历史影响\\
案例实验 & 特定样本在特定任务中的判断分布 & 多数判断自动决定规范真理\\
规范建构 & 研究者提出并辩护的新方案 & 历史作者已经完整提出该方案\\
\bottomrule
\end{tabularx}
\end{center}

\minitopic{“任何比较都使用第三种语言”的反对}最强的方法论反对指出：一旦选取“证据”“能力”或“纠错”作为共同问题轴，比较者已经用自己的元语言重排双方，因此不存在完全中立的比较。这个反对是正确的警告，却不能推出比较不可能。没有中立视角，不等于所有视角同样任意。研究者可以公开桥梁概念的来源，展示它在双方材料中抓住与遗漏什么，并允许对象材料迫使桥梁改写。

反身性因此不是写一段身份声明便完成，而是报告选择怎样影响结果。为什么选这两份文本而不是内部反例？为什么用“知识”而不是“学习”作轴？若改用另一译法，结论哪里变化？当这些问题有可检查答案时，比较仍然能够获得有程度的客观性：结论依赖哪些证据、替代框架会怎样改变它，都向批评者开放。

\minitopic{失败的对应怎样改进理论}假定研究最终发现“庄周梦蝶就是笛卡尔怀疑论”无法成立：叙事角色、论证目标和回应方式都不同。结果不应被压缩成“不可比较”。失败提示我们，原桥梁“梦境使知觉不可靠”过窄；更合适的问题也许是“主体如何理解视角转换以及自己无法占据的立场”。新问题既能保存差异，也可能揭示当代怀疑论忽略的主体转化。

负面比较的价值在于改变问题，而不是勉强保留对应。它需要说明哪项证据推翻了原假设、较弱对应是否仍成立，以及双方理论分别因此看见什么盲点。若每次失败都被改写成更抽象的“其实相似”，理论就无法被反驳；若第一次差异便终止研究，又无法学习。可修订性要求比较者预先说明什么会使自己缩小、改写或放弃结论。

\begin{argumentbox}
一份成熟的比较结论应同时包含五句话：我比较的具体对象是什么；共同问题轴为何成立；最强相似证据是什么；最强差异或竞争解释是什么；这些材料怎样迫使我修改自己的问题。缺少最后一句，比较往往只是在为既有理论寻找异域例证。
\end{argumentbox}

\index{桥梁问题}
\index{竞争解释}
\index{跨文化案例实验}
\index{反身性}
